\section{Conclusion and future work}\label{sec:conclusion}

\pla{} turns a certainty-factor-style calculus into a verifiable
research vehicle: explicit semantics whose formulas execute as tests,
machine-checked convergence, saturation-free log-odds context
conditioning, convex weight learning tied exactly to the engine, and
explanation fidelity measured against a control --- all inside a
dependency-free package whose every reported number regenerates from
committed scripts.

The real-data studies are in (Section~\ref{sec:evaluation}): eleven
readable rules $0.002$--$0.034$ per split below the interpretable
accuracy frontier on credit-card fraud across ten splits, calibration
that beats the constant-prevalence floor, rule-level traces that beat
reversed and random controls on every dataset --- and a pre-registered
negative on the context-conditioning mechanism itself, whose paired
ablation intervals bound any effect to practical irrelevance under all
three designs. The
negative scopes the next iteration precisely: richer context features
than single-bit indicators, rule re-mining under vocabulary drift, and
the full serial-fraud protocol of~\cite{bao2020}. Beyond those, three
directions are specified with hypotheses and kill criteria in the
repository's gap statement: negation in the symbolic layer (making
constraint gating real), a human-subject study of trace utility in
audit decisions, and \pla{} as the uncertain-rule verifier in LLM
generate-and-verify pipelines~\cite{kambhampati2024}.
